% Part: computability
% Chapter: recursive-functions
% Section: notations-pr-functions

\documentclass[../../../include/open-logic-section]{subfiles}

\begin{document}

\olfileid{cmp}{rec}{not}
\olsection{आदिम पुनरावर्ती फलनांची चिन्हांकने}

आदिम पुनरावर्ती फलनांचे अचूक विगमनाधारित वर्णन उपलब्ध असल्याचा एक फायदा
म्हणजे त्यांचे वर्णन पद्धतशीरपणे करता येते. उदाहरणार्थ, पुढीलप्रमाणे
प्रत्येक अशा फलनाला एक ``चिन्हांकन'' देता येते. शून्य, उत्तरवर्ती आणि
प्रक्षेपणे यांसाठी अनुक्रमे $\Zero$, $\Succ$ आणि $\Proj{n}{i}$ ही चिन्हे
वापरा. आता $h$ हे $k$-स्थानी फलन~$f$ आणि $n$-स्थानी फलने $g_0$,
\dots,~$g_{k-1}$ यांच्यापासून संयोजनाने परिभाषित केलेले आहे, आणि या
फलनांना अनुक्रमे $F$, $G_0$, \dots,~$G_{k-1}$ ही चिन्हांकने दिली आहेत
असे समजा. मग नवीन चिन्ह $\fn{Comp}_{k,n}$ वापरून फलन $h$ हे
$\fn{Comp}_{k,n}[F,G_0,\dots,G_{k-1}]$ ने दर्शवता येते.

आदिम पुनरावर्तनाने परिभाषित केलेल्या फलनांसाठी समरूप चिन्हांकने वापरता
येतात. $(k+1)$-स्थानी फलन~$h$ हे $k$-स्थानी फलन~$f$ आणि $(k+2)$-स्थानी
फलन~$g$ यांच्यापासून आदिम पुनरावर्तनाने परिभाषित केलेले आहे, आणि $f$
व~$g$ यांना अनुक्रमे $F$ आणि~$G$ ही चिन्हांकने दिली आहेत असे समजा. मग
$h$ ला दिलेले चिन्हांकन $\fn{Rec}_k[F,G]$ असते.

बेरीज फलनाची आदिम पुनरावर्तनाने व्याख्या पुढीलप्रमाणे आहे, हे आठवा:
\begin{align*}
  \Add(x_0, 0) & = \Proj{1}{0}(x_0) = x_0\\
  \Add(x_0, y+1) & = \Succ(\Proj{3}{2}(x_0, y, \Add(x_0, y))) = \Add(x_0, y) +1
\end{align*}
येथे~$f$ ची भूमिका $\Proj{1}{0}$ बजावते आणि~$g$ ची भूमिका
$\Succ(\Proj{3}{2}(x_0, y, z))$ बजावते. हे $1$-स्थानी फलन~$\Succ$
आणि $3$-स्थानी फलन~$\Proj{3}{2}$ यांच्यापासून संयोजनाने परिभाषित
केलेल्या फलनाचा परिणाम असल्यामुळे त्याला
$\fn{Comp}_{1,3}[\Succ,\Proj{3}{2}]$ हे चिन्हांकन दिले जाते. या
व्यवस्थेत बेरीज फलन पुढील चिन्हांकनाने दर्शवता येते:
\[
\fn{Rec}_1[\Proj{1}{0},\fn{Comp}_{1,3}[\Succ,\Proj{3}{2}]].
\]
अशी चिन्हांकने कधीकधी उपयोगी ठरतात; उदाहरणार्थ, आदिम पुनरावर्ती फलनांचे
प्रगणन करताना.

\begin{prob}
  $\Mult$ साठी संपूर्ण आदिम पुनरावर्ती चिन्हांकन द्या.
\end{prob}

\end{document}
